\documentclass[10pt, letterpaper]{article}

% Packages:
\usepackage[
    ignoreheadfoot, % set margins without considering header and footer
    top=2 cm, % seperation between body and page edge from the top
    bottom=2 cm, % seperation between body and page edge from the bottom
    left=2 cm, % seperation between body and page edge from the left
    right=2 cm, % seperation between body and page edge from the right
    footskip=1.0 cm, % seperation between body and footer
    % showframe % for debugging 
]{geometry} % for adjusting page geometry
\usepackage{titlesec} % for customizing section titles
\usepackage{tabularx} % for making tables with fixed width columns
\usepackage{array} % tabularx requires this
\usepackage[dvipsnames]{xcolor} % for coloring text
\definecolor{primaryColor}{RGB}{0, 79, 144} % define primary color
\usepackage{enumitem} % for customizing lists
\usepackage{fontawesome5} % for using icons
\usepackage{amsmath} % for math
\usepackage[
    pdftitle={Aadesh Tikhe's Master Resume},
    pdfauthor={Aadesh Tikhe},
    pdfcreator={LaTeX with RenderCV},
    colorlinks=true,
    urlcolor=primaryColor
]{hyperref} % for links, metadata and bookmarks
\usepackage[pscoord]{eso-pic} % for floating text on the page
\usepackage{calc} % for calculating lengths
\usepackage{bookmark} % for bookmarks
\usepackage{lastpage} % for getting the total number of pages
\usepackage{changepage} % for one column entries (adjustwidth environment)
\usepackage{paracol} % for two and three column entries
\usepackage{ifthen} % for conditional statements
\usepackage{needspace} % for avoiding page brake right after the section title
\usepackage{iftex} % check if engine is pdflatex, xetex or luatex

% Ensure that generate pdf is machine readable/ATS parsable:
\ifPDFTeX
    \input{glyphtounicode}
    \pdfgentounicode=1
    % \usepackage[T1]{fontenc} % this breaks sb2nov
    \usepackage[utf8]{inputenc}
    \usepackage{lmodern}
\fi



% Some settings:
\AtBeginEnvironment{adjustwidth}{\partopsep0pt} % remove space before adjustwidth environment
\pagestyle{empty} % no header or footer
\setcounter{secnumdepth}{0} % no section numbering
\setlength{\parindent}{0pt} % no indentation
\setlength{\topskip}{0pt} % no top skip
\setlength{\columnsep}{0cm} % set column seperation
\makeatletter
\let\ps@customFooterStyle\ps@plain % Copy the plain style to customFooterStyle
\patchcmd{\ps@customFooterStyle}{\thepage}{
    \color{gray}\textit{\small Aadesh Tikhe - Page \thepage{} of \pageref*{LastPage}}
}{}{} % replace number by desired string
\makeatother
\pagestyle{customFooterStyle}

\titleformat{\section}{\needspace{4\baselineskip}\bfseries\large}{}{0pt}{}[\vspace{1pt}\titlerule]

\titlespacing{\section}{
    % left space:
    -1pt
}{
    % top space:
    0.3 cm
}{
    % bottom space:
    0.2 cm
} % section title spacing

\renewcommand\labelitemi{$\circ$} % custom bullet points
\newenvironment{highlights}{
    \begin{itemize}[
        topsep=0.10 cm,
        parsep=0.10 cm,
        partopsep=0pt,
        itemsep=0pt,
        leftmargin=0.4 cm + 10pt
    ]
}{
    \end{itemize}
} % new environment for highlights

\newenvironment{highlightsforbulletentries}{
    \begin{itemize}[
        topsep=0.10 cm,
        parsep=0.10 cm,
        partopsep=0pt,
        itemsep=0pt,
        leftmargin=10pt
    ]
}{
    \end{itemize}
} % new environment for highlights for bullet entries


\newenvironment{onecolentry}{
    \begin{adjustwidth}{
        0.2 cm + 0.00001 cm
    }{
        0.2 cm + 0.00001 cm
    }
}{
    \end{adjustwidth}
} % new environment for one column entries

\newenvironment{twocolentry}[2][]{
    \onecolentry
    \def\secondColumn{#2}
    \setcolumnwidth{\fill, 4.5 cm}
    \begin{paracol}{2}
}{
    \switchcolumn \raggedleft \secondColumn
    \end{paracol}
    \endonecolentry
} % new environment for two column entries

\newenvironment{header}{
    \setlength{\topsep}{0pt}\par\kern\topsep\centering\linespread{1.5}
}{
    \par\kern\topsep
} % new environment for the header

% Add near your other custom commands (before \begin{document} is fine):
\def\monthname{%
    \ifcase\month
    \or January\or February\or March\or April\or May\or June%
    \or July\or August\or September\or October\or November\or December%
    \fi
}

\newcommand{\placelastupdatedtext}{% \placetextbox{<horizontal pos>}{<vertical pos>}{<stuff>}
  \AddToShipoutPictureFG*{% Add <stuff> to current page foreground
    \put(
        \LenToUnit{\paperwidth-2 cm-0.2 cm+0.05cm},
        \LenToUnit{\paperheight-1.0 cm}
    ){\vtop{{\null}\makebox[0pt][c]{
        \small\color{gray}\textit{Last updated in \monthname\ \the\year}\hspace{\widthof{Last updated in September 2024}}
    }}}%
  }%
}%

% save the original href command in a new command:
\let\hrefWithoutArrow\href

% new command for external links:
\renewcommand{\href}[2]{\hrefWithoutArrow{#1}{\ifthenelse{\equal{#2}{}}{ }{#2 }\raisebox{.15ex}{\footnotesize \faExternalLink*}}}


\begin{document}
    \newcommand{\AND}{\unskip
        \cleaders\copy\ANDbox\hskip\wd\ANDbox
        \ignorespaces
    }
    \newsavebox\ANDbox
    \sbox\ANDbox{}

    \placelastupdatedtext
    \begin{header}
        \textbf{\fontsize{24 pt}{24 pt}\selectfont Aadesh Tikhe}

        \vspace{0.3 cm}

        \normalsize
        \mbox{{\color{black}\footnotesize\faMapMarker*}\hspace*{0.13cm}Chhatrapati Sambhajinagar, India}%
        \kern 0.25 cm%
        \AND%
        \kern 0.25 cm%
        \mbox{\hrefWithoutArrow{mailto:aadeshtikhe24@gmail.com}{\color{black}{\footnotesize\faEnvelope[regular]}\hspace*{0.13cm}aadeshtikhe24@gmail.com}}%
        \kern 0.25 cm%
        \AND%
        \kern 0.25 cm%
        \mbox{\hrefWithoutArrow{tel:+91 9130511165}{\color{black}{\footnotesize\faPhone*}\hspace*{0.13cm}+91 9130511165}}%
        \kern 0.25 cm%
        \AND%
        \kern 0.25 cm%
        \mbox{\hrefWithoutArrow{https://aadesh24.com}{\color{black}{\footnotesize\faGlobe}\hspace*{0.13cm}aadesh24.com}}%
        \kern 0.25 cm%
        \AND%
        \kern 0.25 cm%
        \mbox{\hrefWithoutArrow{https://orcid.org/0009-0009-3185-0845}{\color{black}{\footnotesize\faOrcid}\hspace*{0.13cm}0009-0009-3185-0845}}%
        \kern 0.25 cm%
        \AND%
        \kern 0.25 cm%
        \mbox{\hrefWithoutArrow{https://www.linkedin.com/in/ast12212224}{\color{black}{\footnotesize\faLinkedinIn}\hspace*{0.13cm}Aadesh Tikhe}}%
        \kern 0.25 cm%
        \AND%
        \kern 0.25 cm%
        \mbox{\hrefWithoutArrow{https://github.com/AST12212224}{\color{black}{\footnotesize\faGithub}\hspace*{0.13cm}AST12212224}}%
        \kern 0.25 cm%
        \AND%
        \kern 0.25 cm%
        \mbox{\hrefWithoutArrow{https://medium.com/@aadeshtikhe24}{\color{black}{\footnotesize\faMedium}\hspace*{0.13cm}aadeshtikhe24}}%
        \kern 0.25 cm%
        \AND%
        \kern 0.25 cm% 
        \mbox{\hrefWithoutArrow{https://math.stackexchange.com/users/1380480/aadesh-tikhe}{\color{black}{\footnotesize\faStackExchange}\hspace*{0.13cm}aadesh-tikhe}}%
        
    \end{header}

    \vspace{0.3 cm - 0.3 cm}



\section{About Me!}

        \begin{onecolentry}{
            I'm obsessively curious. I can't see a mathematical pattern without wanting to pull it apart and see if it connects to something else. My work is a mix of rigorous research (published on OEIS, presented at conferences) and experimental coding projects where I just try things. I think with whatever knowledge I have, recombine ideas in weird ways, and occasionally stumble onto something that works. I'm not just interested in solving known problems. I want to find the problems no one's asking yet. My north star is building things that are both mathematically elegant and practically useful.}

            \textbf{Goal:} Advance interpretability in deep learning by applying insights from number theory and algorithmic analysis to neural network behavior. I aspire to build ethical, human-centered systems that simplify complexity and deepen understanding.

            \textbf{Practice:} Committed Vipassana meditator: completed 3 ten-day silent residential courses and served in 1. A discipline that shapes how I approach sustained, deep-focus research.
        \end{onecolentry}
        
\section{Research Highlights}

    \begin{onecolentry}
        \begin{highlights}
            \item \textbf{A consistent method across four independent research lines:} take a classical discrete structure (Pascal's triangle, perfect squares, prime factorization, spatial polytopes) and apply difference or floor/ceiling operations to expose a hidden unifying structure underneath it.
            \item \textbf{Layered Pascal Transform Triangles} unified five previously disconnected OEIS sequences (including one of the author's own closed-form contributions, A259569) under a single generating framework, closing a gap that had gone unnoticed since Voigt (1901) and Piza (1947).
            \item \textbf{DSRS} (floor/ceiling structures on perfect squares) was reviewed by \textbf{Simon Plouffe} (OEIS co-contributor), who called it "not trivial" and forwarded it to \textbf{Jean-Paul Allouche} (CNRS) for further analysis.
            \item \textbf{C2 Conjecture}, a recursive prime decomposition result verified across 1 billion+ primes with zero counterexamples, drew sustained engagement from \textbf{Paul Lockhart} (\textit{A Mathematician's Lament}), who assessed it as comparably difficult to open problems like $p = n^2+1$.
            \item \textbf{AIDEAS-OEIS} extends this same instinct (pattern-checking and structural validation) into a working AI system, applying LLM + symbolic verification to flag errors across 2,000+ OEIS entries.
        \end{highlights}
    \end{onecolentry}
    
\section{Entrepreneurship}

\begin{twocolentry}{
    \textit{July 2026 – Present}}
    \textbf{Founder \& Managing Director, ATAEON Cognition Private Limited} \\
    \textit{DPIIT-recognized startup (Startup India), Certificate No. DIPP282717}
\end{twocolentry}

\vspace{0.10 cm}

\begin{onecolentry}
    \begin{highlights}
        \item Founded and incorporated ATAEON Cognition Private Limited (20 July 2026); recognized as a startup by the \textbf{Department for Promotion of Industry and Internal Trade (DPIIT)}, Government of India, on 13 September 2026, in the Human Resources industry, Recruitment sector. \href{https://aadesh24.com/certificates/dpiit-startup-recognition-2026.pdf}{Certificate Link}
        \item Built \textbf{LabourVala} (\href{https://labourvala.com}{labourvala.com}), the company's first product: an ethical hiring platform connecting skilled daily-wage workers with verified contractors in Chhatrapati Sambhajinagar. The product is complete and in its launch phase
        \item Contractors and companies are verified with licences and documents, and workers with Aadhaar and PAN; phone numbers are revealed only after both sides accept, which keeps matches on the platform
        \item 4-digit arrival codes create a digital attendance record, two-way ratings build trust scores, and a flat fee is charged only per successful match; Marathi-first web and Android apps
    \end{highlights}
\end{onecolentry}

    \section{Education \& Academic Achievements}



        
        \begin{twocolentry}{
            \textit{Aug 2023 – June 2026}}
            \textbf{Amity University, Mumbai} \\
            \textit{Bachelor of Computer Applications (BCA)}
        \end{twocolentry}
        
        \vspace{0.1cm}
        
        \begin{onecolentry}
            \begin{highlights}
                \item CGPA: 8.39 / 10.0 (till date)
                \item \textbf{Key Coursework:} Operating Systems, Theory of Computation, Comparative Learning Algorithms, Linear Algebra for ML
            \end{highlights}
        \end{onecolentry}

        \vspace{0.2cm}
        
        \begin{twocolentry}{
            \textit{May 2017 – May 2022}}
            \textbf{Wockhardt Global School (IB), Chhatrapati Sambhajinagar} \\
            \textit{IB Diploma Programme (DP), 2020–2022}: Final Grade 30 / 45 \\
            \textit{IB Middle Years Programme (MYP), 2017–2020}: Final Grade 85\% \\
            \textbf{Head Boy \& Leadership Roles}
        \end{twocolentry}
        
        \vspace{0.1cm}
        
        \begin{onecolentry}
            \begin{highlights}
                \item \textbf{Extended Essay:} Analyzed the time complexity of the NegaMax algorithm (Strazilla) vs. Alpha-Beta Pruning (NNUE Stockfish)
                \item \textbf{Main Subjects:} Computer Science HL, Mathematics: AA HL, Business Management HL
                \item Led school initiatives as \textbf{Head Boy}, including playground renovation and MUN 2020 organization
                \item Represented school at \textbf{Harvard Model United Nations (HMUN) 2018}, an international MUN conference
                \item \textbf{Chaired UNICEF at WGS MUN} 2021; facilitated engaging agendas and resolutions
                \item Served on the \textbf{Board of Directors, Rotaract Club Aurangabad}, fostering leadership and collaboration
                \item Participated in the \textbf{Youth India Express Summit}, a national-level youth leadership event
                \item Handled \textbf{stage management and decoration at TEDxYouth@WockhardtGlobalSchool} (Nov 2021); oversaw technical execution, speaker coordination and stage design
                \item Initiated \textbf{Stevia The Developer} (2019–2020, MYP 5 / Grade 10): Marathwada faces chronic water scarcity, and sugarcane, its main cash crop, consumes large amounts of water. Researched Stevia as a low-water natural sweetener, grew and tested it on our own farmland, and ran awareness sessions for farmers on its water, health and economic benefits; overcame stage fear by addressing these diverse audiences
            \end{highlights} 
        \end{onecolentry}


\section{Curiosity-Driven Research \& Publications}
        
\begin{twocolentry}{
    \textit{Aug 2022 - Dec 2024}
}
    \textbf{Probability in Regular 2-Polytopes}

    \vspace{0.10 cm}
\end{twocolentry}

\vspace{0.10 cm}

\begin{onecolentry}
    \begin{highlights}
        \item Independently conducted research from 2022–2024 on a novel spatial probability model within regular 2D polytopes
        \item Proposed a Periodic Cotangent Function to model the probability distribution of a 0 polytope relative to the centroid and boundary
        \item Presented this work at the \textbf{Indian Mathematical Society Conference} (MIT-WPU, Pune) on \textbf{25 December 2024}, gaining valuable feedback and insights
        \item Full explanation available via presentation video: \href{https://youtu.be/Vwe1ojJnU_A?si=o3FbliMeOJcqa_1W}{youtu.be/Vwe1ojJnU\_A}; written paper still in progress
    \end{highlights}
\end{onecolentry}

\vspace{0.1cm}

\begin{twocolentry}{
    \textit{Feb 2024 – 2025}}
    \textbf{Layered Pascal Transform Triangles: A Unified Framework for Iterated Finite Differences on Monomials}
\end{twocolentry}

\vspace{0.10 cm}

\begin{onecolentry}
    \begin{highlights}
        \item Organized iterated forward differences $\Delta^\ell(x^p)$ as a parametrized family of triangular coefficient arrays indexed by layer $\ell$, where Layer 0 recovers classical Pascal's triangle and higher layers yield transformed structures retaining Pascal-like organization
        \item Derived a closed-form coefficient generator $C_{p,\ell}(v)$ that directly produces individual polynomial coefficients without requiring polynomial expansion or Stirling number basis conversion, enabling $O(\ell)$ access to specific coefficients when $p \gg \ell$
        \item Unified five previously disconnected OEIS sequences (A259569, the author's own closed-form contribution, plus A001117, A000918, A000919 and A000920) as fixed-layer projections of a single underlying structure
        \item Connected the framework to historical work by Voigt (1901, ``Konformanten'') and Piza (1947, ``D-coefficients''), showing both contain numerical values matching Layers 2 and 3 but lacked systematic organization, direct coefficient generation, or unification across sequences
        \item Entire construction uses only elementary combinatorial tools (binomial coefficients, forward difference operator), accessible at the undergraduate level while providing computational utility for sparse coefficient extraction
        \item Originated from informal exploration of Pascal's triangle patterns in 2023; developed into a formal framework starting February 2024
        \item Preprint archived on Zenodo with DOI: \href{https://doi.org/10.5281/zenodo.18256084}{10.5281/zenodo.18256084}
    \end{highlights}
\end{onecolentry}

\begin{twocolentry}{
    \textit{Apr 2025 – Present}
}
    \textbf{Fibonacci Numbers from Pascal Rows: A Ternary Coefficient Approach}
\end{twocolentry}

\vspace{0.10 cm}

\begin{onecolentry}
    \begin{highlights}
        \item Established novel algebraic framework for transforming Pascal triangle rows into Fibonacci numbers through systematic coefficient selection where $c_i \in \{-1, 0, 1\}$
        \item Developed six complete enumeration strategies with CUDA GPU acceleration, achieving exhaustive analysis up to computational limit $3^{22}$ (31.4 billion combinations)
        \item Designed a greedy algorithm leveraging palindromic symmetry and center-dominance properties as a candidate $O(n)$ heuristic for the exponential $O(3^n)$ search space
        \item Later identified cases where the greedy heuristic does not reliably reproduce correct coefficient selections; currently investigating the boundary conditions under which it holds and does not hold
        \item Preprint archived on Zenodo with DOI: \href{https://doi.org/10.5281/zenodo.17412193}{10.5281/zenodo.17412193}
        \item Code implementation available at \href{https://github.com/AST12212224/Pascals-triangle-row-for-Fibonacci}{GitHub Repository}
    \end{highlights}
\end{onecolentry}

\begin{twocolentry}{
    \textit{June 2025 – Sept 2025}
}

    \textbf{Discrete Square Residual Structures (DSRS): A Framework Where Every Integer Reveals Its Own Connection to $\pi$}
    
    \vspace{0.10 cm}
\end{twocolentry}

\vspace{0.10 cm}

\begin{onecolentry}
    \begin{highlights}
        \item Introduced a purely arithmetic framework using floor and ceiling operations on perfect squares: for fixed $\mu$, defined $U(n) = \lceil n^2/\mu \rceil$ and $L(n) = \lfloor n^2/\mu \rfloor$, producing infinite products $P(\mu) = \prod_{n=n_0}^{\infty} \frac{\Delta_L(n)}{\Delta_U(n)}$
        \item Computational evidence up to $N = 10^6$ and $\mu \leq 70{,}000$ reveals a sharp dichotomy: $P(\mu) \to \pi/2$ when $\mu \not\equiv 0 \pmod{4}$, and $P(\mu) \to 1$ when $\mu \equiv 0 \pmod{4}$ or $\mu = 1$
        \item Provided structural argument showing all non-multiples of 4 reduce to the $\mu = 2$ case (a floor-ceiling encoding of the Wallis product) via residual cancellation, with increasing entropy for larger $\mu$ requiring more terms to stabilize
        \item Identified OEIS correspondences: $\mu = 2$ produces A052928, A007590, A000982, A063196; $\mu = 3$ yields A004523, A000212, A008810, A130823; $\mu = 4$ yields A008619, A002620, A004652, A161239
        \item Research reviewed by \textbf{Simon Plouffe} (OEIS co-contributor, Inverse Symbolic Calculator), who described the work as ``an interesting topic'' and ``not trivial,'' forwarded the paper to \textbf{Jean-Paul Allouche} (CNRS, expert in infinite products and automatic sequences) for further analysis, and offered an ArXiv recommendation
        \item Preprint archived on Zenodo with DOI: \href{https://doi.org/10.5281/zenodo.17519713}{10.5281/zenodo.17519713}; complete source code at \href{https://github.com/AST12212224/DSRS}{GitHub Repository}
    \end{highlights}
\end{onecolentry}


\begin{twocolentry}{
    \textit{May 2026 – Present}}
    \textbf{A Recursive Prime Decomposition Conjecture (C2)}
\end{twocolentry}

\vspace{0.10 cm}

\begin{onecolentry}
    \begin{highlights}
        \item Proposed Conjecture C2: for every prime $p > 2$, there exist $q, r \in \{1\} \cup \{\text{primes}\}$ with $q < p$ such that $p = 2q + r$ or $p = 2q - r$; since $r$ also decomposes, every prime unfolds into a binary tree terminating at leaves of 1
        \item Formulated the strictly stronger C2-Additive variant: additive only ($p = 2q + r$, $q < p/2$), eliminating $\pm$ ambiguity and producing canonical single-direction decomposition trees
        \item Conjecture is strictly stronger than Lemoine's conjecture (1895) restricted to the prime domain: C2 requires $q < p$, Lemoine does not; if C2 is proved, Lemoine for primes follows as a corollary
        \item Verified C2 for all 1,008,309,544 primes (up to $\sim 2 \times 10^{10}$) with \textbf{zero counterexamples}, using a two-stage CUDA pipeline on NVIDIA RTX 5070: sliding-window search (500 candidates) followed by brute-force on failures with Miller-Rabin primality testing, achieving $\sim$150M primes/sec throughput; built interactive HTML visualizations of recursive decomposition trees for exploratory analysis
        \item Reviewed by \textbf{Paul Lockhart} (mathematician, author of \textit{A Mathematician's Lament}), who assessed the conjecture as ``possibly even more difficult'' than classical Diophantine problems like $p = n^2 + 1$, situating it within the structural gap between additive and multiplicative number theory
        \item Posed the open question: does iterative generation from $S_0 = \{1, 2\}$ via $2q + r$ produce exactly the primes? If yes, this would yield a structural definition of primality independent of the Fundamental Theorem of Arithmetic
        \item Preprint archived on Zenodo with DOI: \href{https://doi.org/10.5281/zenodo.20225703}{10.5281/zenodo.20225703}; source code at \href{https://github.com/AST12212224}{GitHub}
    \end{highlights}
\end{onecolentry}


\begin{twocolentry}{
    \textit{In progress}}
    \textbf{Dimensional Primality: A Geometric Hierarchy of Integers via Discrete Hyper-Rectangles}
\end{twocolentry}

\vspace{0.10 cm}

\begin{onecolentry}
    \begin{highlights}
        \item Interprets integer factorization geometrically: $N$ as the volume of a discrete hyper-rectangle with all integer sides greater than 1; proved the maximum dimension in which $N$ admits such a genuine representation is exactly $\Omega(N)$, the number of prime factors counted with multiplicity
        \item Defined a strictly nested hierarchy $D_1 \supset D_2 \supset D_3 \supset \cdots$ with $D_n = \{N : \Omega(N) \geq n\}$, and \textit{$n$-dimensional primes} as the integers outside $D_n$ (1 together with the $k$-almost primes, $k \leq n-1$), recovering the classical primes at $n = 2$
        \item Proved a forced-collapse corollary (any representation in more than $\Omega(N)$ dimensions has at least $n - \Omega(N)$ unit sides) and that the smallest member of $D_n$ is $2^n$, attained uniquely by the cube $(2, \ldots, 2)$
        \item Working paper archived on Zenodo with DOI: \href{https://doi.org/10.5281/zenodo.20645595}{10.5281/zenodo.20645595}
    \end{highlights}
\end{onecolentry}

\section{Blog \& Expository Writing}

\begin{twocolentry}{
    \textit{Dec 2024 – Apr 2025}
}
    \textbf{Efficient Fibonacci Computation: Édouard Lucas's Pascal Triangle Method} (Medium)
\end{twocolentry}

\vspace{0.10 cm}

\begin{onecolentry}
    \begin{highlights}
        \item Analyzed Lucas's combinatorial approach for calculating the exact $n^{th}$ Fibonacci number without recursion, comparing computational complexity against classical methods (matrix exponentiation, Binet's formula, naive recursion)
        \item Demonstrated feasibility in resource-constrained environments; code available at \href{https://github.com/AST12212224/Efficient-Fibonacci-N-th-term-calculation-E_Lucas.git}{GitHub Repository}
        \item Full write-up: \href{https://medium.com/@aadeshtikhe24/efficient-fibonacci-computation-instantly-calculate-the-exact-nth-number-directly-no-recursion-ed13bb284c71}{Medium Blog}
    \end{highlights}
\end{onecolentry}

\vspace{0.1cm}

\begin{twocolentry}{
    \textit{Dec 2025}
}
    \textbf{Escaping the Killer by Always Getting Closer} (Self-published)
\end{twocolentry}

\vspace{0.10 cm}

\begin{onecolentry}
    \begin{highlights}
        \item Expository paper exploring how a locally ``greedy'' rule (always decrease distance to target) can produce paths of arbitrarily large, even infinite, total length, using the Archimedean spiral as the constructive example
        \item Derived the closed-form arc length formula $L = \frac{\pi R^2}{m}\sqrt{1 + \left(\frac{m}{2\pi R}\right)^2} + \frac{m}{4\pi}\ln\!\left(\frac{2\pi R}{m}\left(1 + \sqrt{1 + \left(\frac{m}{2\pi R}\right)^2}\right)\right)$ and showed $L \to \infty$ as the radial step $m \to 0^+$
        \item Connected the result to gradient descent pathologies: locally improving loss at every step does not bound total path complexity or guarantee convergence in practical time
        \item Available at \href{https://aadesh24.com/Escaping_the_Killer_by_Always_Getting_Closer.pdf}{aadesh24.com}
    \end{highlights}
\end{onecolentry}

    
\section{Projects and Contributions}

\begin{twocolentry}{
    \textit{Jan 2026 – May 2026}
}
    \textbf{AIDEAS-OEIS: AI-Driven Evaluation of Arithmetic Sequences in OEIS} (Python, LLMs)
\end{twocolentry}

\vspace{0.10 cm}

\begin{onecolentry}
    \begin{highlights}
        \item Developed a hybrid verification system combining \textbf{Qwen2.5-Math:7b} (served locally via Ollama) with SymPy symbolic computation and LODA program synthesis to automatically interpret, validate, and flag erroneous formulas across OEIS entries
        \item Built a three-verifier pipeline: LODA algorithmic verification ($\sim$19\% hit rate) $\to$ LLM formula-to-code generation in Restricted Python sandbox $\to$ SymPy closed-form evaluation, reconciled via Bayesian classifier with confidence scores
        \item Validated 2,041 OEIS sequences: \textbf{51.6\%} HIGH\_CONFIDENCE\_VALID, 40.3\% NEEDS\_HUMAN\_REVIEW, 8.1\% LLM\_INTERPRETATION\_ERROR, with Wilson 95\% confidence intervals confirming statistical stability
        \item Designed with crash-resume safety (SQLite persistence per sequence), Docker containerization, and CSV/JSON /SQLite export for community-scale deployment
        \item BCA final-year capstone project (7 credits, scored 90/100); open-source at \href{https://github.com/AST12212224/AIDEAS-OEIS}{GitHub Repository}
    \end{highlights}
\end{onecolentry}


        \begin{twocolentry}{
            \textit{\href{https://oeis.org/wiki/User:Aadesh_Tikhe}{oeis.org/wiki/User:Aadesh}}}
            \textbf{OEIS Contributor – Integer Sequence Research}
        \end{twocolentry}
        
        \begin{onecolentry}
            \begin{highlights}
                \item \textbf{Authored closed-form formula for OEIS} \href{https://oeis.org/A259569}{A259569}; later shown to be a Layer 2 projection of the author's unified Pascal transform framework (see Layered Pascal Transform Triangles above)
                \item \textbf{Corrected mathematical inaccuracies in} \href{https://oeis.org/A130823}{A130823}, improving sequence integrity and documentation
                \item Active contributor on the \href{https://oeis.org/wiki/User:Aadesh_Tikhe}{OEIS Wiki}, a globally recognized encyclopedia of integer sequences
            \end{highlights}
        \end{onecolentry}

        \vspace{0.1cm}

        \begin{twocolentry}{
            \textit{Ongoing}
        }
            \textbf{Manim Codes for Math Research Presentation }(Python [Manim], C++)
        \end{twocolentry}
        
        \vspace{0.10 cm}
        
        \begin{onecolentry}
            \begin{highlights}
                \item Built 10+ Manim animations to visualize novel mathematical concepts (Fibonacci methods, probability distributions), garnering 500+ views on YouTube.
                \item Used in YouTube presentation: \href{https://youtu.be/Vwe1ojJnU_A?si=0aWaau3x-QwqNVsj}{Watch here}.
                \item \href{https://github.com/AST12212224/Manim-codes-for-my-maths-research-presentation}{GitHub Repository}
            \end{highlights}
        \end{onecolentry}

        \begin{twocolentry}{
            \textit{2021 – 2022}
        }
            \textbf{Django-POL\_Agro Web Platform }(Python, HTML, CSS)
        \end{twocolentry}
        
        \vspace{0.10 cm}
        
        \begin{onecolentry}
            \begin{highlights}
                \item Built for a client during IBDP; developed a full-stack Django-based CRUD platform.
                \item Connects fertilizer wholesalers with farmers to streamline queries and product access.
                \item Hosted on localhost for demo; source code: \href{https://github.com/AST12212224/Django-POL_Agro}{GitHub Repository}.
            \end{highlights}
        \end{onecolentry}

        
        \begin{twocolentry}{
            \textit{Ongoing}
        }
            \textbf{myBash – Bash \& Linux Commands Repository }(Shell)
        \end{twocolentry}
        
        \vspace{0.10 cm}
        
        \begin{onecolentry}
            \begin{highlights}
                \item A collection of useful Bash commands, Linux one-liners, and \texttt{.bashrc} customizations. Includes a cron job for auto-updating Homebrew monthly. Personal notes and hands-on summaries from tutorials.
                \item \href{https://github.com/AST12212224/myBash}{GitHub Repository}
            \end{highlights}
        \end{onecolentry}
    
    \section{Key Skills}
        
        \begin{onecolentry}
            \textbf{Languages:} C++, C, Python, SQL, Bash, Cuda, Rust \\
            \textbf{Libraries \& Frameworks:} Django, Manim, Pandas, NumPy, Matplotlib, Seaborn, Scikit-learn, Pygame \\
            \textbf{Tools \& Platforms:} Git, Docker, AWS, Google Colab, LaTeX.\\
            \textbf{Core Strengths:} Mathematics, Data Structures \& Algorithms, Artificial Intelligence, Machine Learning, C Programming, Operating Systems (incl. Bash Shell), and Computational Theory \\
            \textbf{Hobbies \& Interests} 
            \begin{itemize}
                \item Painting – Creative expression through various mediums.
                \item Cycling \& Hiking – Fitness and nature exploration.
                \item Piano – Self-taught; composed and performed a piano cover, published on \href{https://youtu.be/MvuJxi9bqvk?si=OAYBCSiWyTuYn6db}{YouTube}.
                \item Meditation – Vipassana practice for mindfulness.
                \item Mathematical Puzzles – Exploring number theory and logic.
            \end{itemize}
            \textbf{Spoken Languages:} English, Hindi, Marathi, German (A1 level) 
        \end{onecolentry}


        
\section{Specialised Courses \& Certifications}

    \begin{onecolentry}
        \begin{itemize}
        \item \textbf{Discrete Mathematics – NPTEL} (IIT Madras, Jan–Apr 2026): 12-week course; scored \textbf{91\%} (Elite certification; one of 979 candidates certified). Topics include combinatorics, graph theory, recurrence relations, and formal logic. \href{https://nptel.ac.in/noc/E_Certificate/NOC26CS68S25130047104321322}{Certificate Link}
        \item \textbf{Theory of Computation – NPTEL} (IIT Hyderabad, Jan–Apr 2026): 12-week course; scored 60\% (Elite certification; one of 262 candidates certified). Topics include automata theory, context-free grammars, Turing machines, and decidability. \href{https://nptel.ac.in/noc/E_Certificate/NOC26CS85S35130004104321322}{Certificate Link}
        \item \textbf{Machine Learning Specialization – Andrew Ng (Stanford Online, DeepLearning.AI)} (Jan 2026): Completed all three courses: supervised learning, advanced algorithms, and unsupervised learning including recommender systems and reinforcement learning. \href{https://coursera.org/share/5b3c25baf30c71e878a3dc5d4d36688a}{Certificate Link}
        \item \textbf{Certified Lean Six Sigma (White Belt)} (Six Sigma Society, July 2023, Cert. No. 76953176): Gained process improvement experience using Minitab, Excel, Python, Lucidchart, PowerPoint, and Microsoft Project. \href{https://aadesh24.com/certificates/lean-six-sigma-white-belt-2023.pdf}{Certificate Link}
        \item \textbf{Certificate of Appreciation – TEDxYouth@WockhardtGlobalSchool} (Nov 2021): for stage management and decoration. \href{https://aadesh24.com/certificates/tedxyouth-wgs-2021.png}{Certificate Link}
        \item All certificates, with verification links: \href{https://aadesh24.com/certifications.html}{aadesh24.com/certifications}
            

        \end{itemize}
    \end{onecolentry}


    

\end{document}